Coordinated information operations increasingly rely on networks of accounts that produce, amplify, and adapt information across online platforms~\cite{cima2024coordinated,graham2024toolkit,minici2024crossplatform}. Large language models (LLMs) further reduce the cost of generating diverse yet persona-consistent content, sustaining interactions, and adjusting behavior in response to platform feedback~\cite{piao2025social,jin2025synthetic}. Consequently, an individual account may appear locally plausible while its relations and activities across multiple accounts over time may reveal patterns of purposeful coordination. This shift makes account-level signals alone an increasingly weak basis for investigating networked activity~\cite{cima2024coordinated,zouzou2024fakefollowers,burghardt2024sociolinguistic}. 

Online coordination is not sufficient evidence of a policy violation~\cite{graham2024toolkit,rogers2025manufactured}. Legitimate organizations, advocacy groups, and public-interest communities may also coordinate online~\cite{graham2024toolkit}. A suspicious operation may combine human-operated accounts, conventional bots, and LLM-driven agents~\cite{zouzou2024fakefollowers,burghardt2024sociolinguistic,piao2025social,jin2025synthetic}. Neither coordination nor automation is therefore conclusive evidence of harmful intent or deceptive identity~\cite{rogers2025manufactured}. We use \emph{suspected coordinated information operations} (\emph{suspected CIOs}) to describe cases in which multiple accounts exhibit potentially purposeful coordination across content, relationships, and time. When additional evidence establishes deceptive identities, concealed orchestration, or relevant platform enforcement, such activity may be classified as coordinated inauthentic behavior (CIB)~\cite{rogers2025manufactured,luceri2026tiktok}. HyperTrace does not infer malicious intent or account authenticity from coordination alone. Instead, it supports professional reviewers in examining the evidence needed for a context-dependent determination.

Professional review of suspected CIOs requires more than detecting coordination. Reviewers should evaluate the evidence that gives the observed coordination its meaning. Such evidence is typically distributed across accounts and interactions~\cite{cima2024coordinated,graham2024toolkit,minici2024crossplatform}. An isolated repost, mention, or shared link may be benign, whereas synchronized actions, repeated interactions, and cross-account diffusion within a particular time window may jointly support an operation-level assessment~\cite{cima2024coordinated,graham2024toolkit,zouzou2024fakefollowers}. Professional reviewers must therefore determine which relationships constitute a relevant coordination mechanism, whether each event was available before the decision cutoff, and whether every evidentiary claim can be verified against an underlying record. % These requirements are derived from our formative interviews and formalized as HyperTrace design requirements.
A risk score or a visually salient subgraph may direct attention to a case, but it cannot establish these evidentiary conditions on its own~\cite{agarwal2023evaluating,schemmer2023appropriate}.

Our formative study confirmed that existing review support may not meet these evidentiary needs. We analyzed structured responses from 10 participants, whose research experience averaged approximately two years. They reported that risk scores and content similarity were insufficient for a final determination. Participants needed to reconstruct how coordination developed over time, inspect the source records underlying individual claims, consider alternative explanations for apparently synchronized behavior, and understand how an assessment changed as new evidence arrived. These findings indicate that review interfaces should support mechanism-based, temporally bounded, provenance-verifiable, and revisable judgments, rather than merely communicate model confidence.

Existing graph explanations may not satisfy the evidentiary requirements of professional review. Prediction-oriented methods can select nodes, edges, or motifs that preserve a model output~\cite{agarwal2023evaluating,chen2023tempme}. Yet an explanation that is sufficient for a prediction may still omit the mechanism by which accounts coordinated, include events that occurred after the decision cutoff, become stale as an operation evolves, or lack links to verifiable source records~\cite{agarwal2023evaluating,wang2024humanexplain}. In high-stakes review, these omissions are not only technical shortcomings. Explanations can make an incorrect recommendation appear persuasive and thereby encourage inappropriate reliance on AI assistance~\cite{schemmer2023appropriate,he2023knowing,schoeffer2024explanations}. An explanation for suspected CIOs should therefore do more than reproduce a detector's output. It should provide evidence that a reviewer can challenge, trace, and independently assess.

The structure of coordinated networks raises a separate question of geometric fidelity. Branching diffusion paths, nested communities, and heterogeneous relations can produce hierarchical and multiscale structures~\cite{choudhary2023hyperbolic,yang2024hypformer,park2024hyperbolic,park2024multihyperbolic}. Hyperbolic representations have been used to model such structures in graph learning and social-bot detection~\cite{choudhary2023hyperbolic,lu2026sahg}. Existing studies have evaluated these representations mainly through predictive performance~\cite{choudhary2023hyperbolic,yang2024hypformer,lu2026sahg}. In this paper, geometric fidelity characterizes the correspondence between the full graph and the evidence subgraph. Visual coherence similarly concerns how an explanation appears to a reviewer rather than how closely it preserves the detector's geometry~\cite{wang2024humanexplain}. We therefore evaluate preservation of the detector's Lorentz geometry as an auditable property of the reconstructed evidence. Our claim concerns explanation fidelity. The comparative classification performance of Lorentz geometry remains an empirical question.

An auditable explanation workflow for suspected CIOs must address the evidentiary and geometric gaps together. Such a workflow must provide compact evidence that preserves the model-relevant geometry, respects the decision-time boundary, and exposes its evidence provenance. These properties are mutually dependent. Geometry without temporal validity may rationalize a decision using future information. Temporal validity without evidence provenance may leave claims unverifiable. Complete evidence without sufficient compression may overwhelm reviewers with the original network. An auditable workflow must therefore satisfy all three requirements simultaneously.

HyperTrace converts a model alert into a compact evidence package that professional reviewers can inspect. The workflow begins with a Lorentz-HGT coordination detector, which models heterogeneous account and content relations in Lorentz space to produce an initial assessment~\cite{hu2020hgt,yang2024hypformer,park2024hyperbolic,lu2026sahg}. HyperTrace then reconstructs a relational snapshot containing the accounts, interactions, and events that support that assessment. The snapshot is \emph{cutoff-safe}: it includes only records available at or before the decision cutoff and excludes later information. A constrained selection procedure reduces the snapshot to a compact evidence subgraph while preserving the detector's assessment and Lorentz decision geometry. Every selected item links to a verifiable source record, making its evidence provenance directly inspectable. The interface presents the detector's output and reconstructed evidence as a preliminary assessment rather than a final verdict. Professional reviewers can therefore inspect the reconstructed mechanism, verify source records, qualify the assessment, and accept or reject the model recommendation.

We ask the following research questions:

\noindent\textbf{RQ1 (Review requirements).}
What evidence and interaction capabilities do professional reviewers need to independently assess suspected CIOs?

\noindent\textbf{RQ2 (Technical fidelity).}
Can constrained evidence reconstruction produce compact evidence subgraphs that preserve Lorentz decision geometry while satisfying temporal-validity and evidence-provenance requirements?

\noindent\textbf{RQ3 (Human decision outcomes).}
Compared with risk scores and conventional static explanations, how does HyperTrace affect reviewers' final decision accuracy, rejection of incorrect model recommendations, review time, and workload?

We evaluated the three research questions through complementary methods. The formative interviews addressed RQ1. For RQ2, we evaluated HyperTrace on controlled LLM-driven coordination simulations and previously unseen operations from real-platform data, keeping entire operations disjoint between development and evaluation. The final offline evaluation comprised 24 high-risk cases from two unseen operations with balanced operation coverage. We measured evidence retention, prediction sufficiency, comprehensiveness, geometric fidelity, temporal validity, provenance coverage, and explanation stability. This evaluation tested whether compact evidence could preserve the required audit properties across both controlled and real-platform settings.

We addressed RQ3 through a controlled review study with 20 domain professionals serving as professional reviewers. The same 24 high-risk cases formed the case pool, and balanced block randomization assigned eight complete cases to each participant. Participants first made an independent judgment without model assistance, then examined the assigned assistance, and finally resubmitted their judgment and confidence. This staged procedure separated initial judgment, evidence-assisted assessment, and final determination. It therefore allowed us to measure both acceptance of correct recommendations and rejection of incorrect ones.

HyperTrace retained a median of only 1.56\% of candidate evidence while preserving the detector's assessment, achieved 0.984 Lorentz geometric fidelity, and provided 100\% provenance coverage. Compared with assistance that provided only a risk score or a conventional explanation, HyperTrace increased professional reviewers' final decision accuracy by 23.9 percentage points and their rejection of incorrect model recommendations by 45 percentage points. These results suggest that evidence compression can satisfy geometric, temporal, and provenance constraints. Making these properties visible may help reviewers evaluate AI recommendations independently and reject incorrect ones.

The main contributions of this paper are as follows:

\begin{itemize}
    \item \textbf{Formative evidence about professional review of suspected CIOs.} We identify four recurring requirements for independent review: mechanism-level evidence, temporal reconstruction, evidence-provenance verification, and revisable assessments.

    \item \textbf{The HyperTrace workflow.} We connect a Lorentz-HGT coordination detector to compact evidence reconstruction under explicit temporal-validity and evidence-provenance constraints, while treating preservation of the detector's Lorentz geometry as an auditable property.

    \item \textbf{A joint technical and human evaluation of appropriate reliance.} We relate evidence-level properties, including prediction sufficiency and geometric fidelity, to reviewers' final accuracy, acceptance of correct recommendations, rejection of incorrect recommendations, review time, and workload under controlled model assistance.
\end{itemize}

High-stakes AI-assisted review requires an evidence-centered standard for explainability. An effective explanation should do more than make a model output appear convincing or visually intelligible. It should enable professional reviewers to verify the evidence, recognize an invalid recommendation, and remain responsible for the final determination.
